% Generated by code/integrate.py; edit extraction rules there.
% Reviewed source: revisions/2026-10-04-r11/manuscript/retained_proofs.tex, line 2, commit f5021cefa71492babfcbfa580e0c984f59a9de26
\section{Proofs}\label{app:proofs}
\subsubsection{Proof of Theorem~\ref{thm:verification}}
Extend the aggregator Lipschitz-continuously outside its invariant interval for this comparison argument. The extension leaves all policy utilities and the candidate residual unchanged. For a nonnegative number $c$, define $d_c(t)=Be^{L(T-t)}+cA_L(T-t)$. Then $d_c(T)=B$, $d_c(t)\geq B$, and $-d_c'=Ld_c+c$. Since a spatially constant shift leaves the generator unchanged and $f$ is $L$-Lipschitz in utility,
\[
 F[v+d_c]\geq F[v]-d_c'-Ld_c=F[v]+c,
 \qquad F[v-d_c]\leq F[v]+d_c'+Ld_c=F[v]-c.
\]
Because $F[v]=r^a[v]-g^a[v]\in[-e-\eta,e]$, $v+d_{e+\eta}$ is a supersolution and $v-d_{e+\eta}$ a subsolution. They dominate the required stopping data on every boundary face. A spatially constant shift also preserves an exactly satisfied homogeneous reflected boundary condition. Comparison proves \eqref{eq:valuebound}.

Apply the same argument to the fixed-policy operator $F^a[v]=-v_t-\HH^a[v]$, with $c=e$, to obtain \eqref{eq:policybound}. The triangle inequality gives \eqref{eq:regret}, and policy admissibility gives its nonnegative lower bound. The construction uses $L\geq0$ and $d_c(t)\geq B$ specifically to cover lateral stopping boundaries; it does not incorrectly discount their errors as if every path stopped only at $T$. For nonsmooth candidates the same argument applies when their residual inequalities are verified in the viscosity sense. The theorem's differential implementation, however, requires the stated smooth candidate.\qed

\subsubsection{Exact policy improvement}
Suppose $v=V^a$ solves the fixed-policy problem and a feasible $\tilde a$ satisfies $\HH^{\tilde a}[v]\geq\HH^a[v]$ pointwise. Then $F^{\tilde a}[v]\leq0$ with the same boundary payoff. Policy comparison gives $V^{\tilde a}\geq V^a$. This is a policy-space statement. A single gradient step on a domain-averaged actor loss does not automatically meet its pointwise premise. The measured action gap and policy-evaluation errors are retained precisely because approximate neural improvement can violate that premise.

\subsubsection{Proof of Theorem~\ref{thm:discrete}}
For a fixed state and action, let $z$ and $\tilde z$ solve \eqref{eq:implicit} for $w$ and $\tilde w$. Strict monotonicity of $z-hf(z)$ gives
\[
 |z-\tilde z|\leq\frac{|P_h^a(w-\tilde w)|}{1+\lambda h}
 \leq q_h\|w-\tilde w\|_\infty.
\]
The same monotonicity implies order preservation. Taking the maximum over actions preserves both order and the Lipschitz constant. Banach's fixed-point argument supplies the two fixed points on the invariant complete value domain.

Now $\|v-v_h^a\|\leq\|v-T_h^av\|+q_h\|v-v_h^a\|$, which gives the first bound. Also $\|v-T_hv\|\leq\delta+\epsilon$, giving the second. The triangle inequality gives the third. All norms are over every state of the specified finite model. An omitted action changes $T_h$ and must be accounted for as a separate error, not hidden in $\delta$.\qed

\subsubsection{Proof of Theorem~\ref{thm:viscosity}}
First, \eqref{eq:discretebound} gives
\[
 \|v_h-v_h^*\|_\infty\leq\frac{\delta_h+\epsilon_h}{1-q_h}
 =\frac{1+\lambda h}{\lambda}\frac{\delta_h+\epsilon_h}{h}\longrightarrow0.
\]
It remains to identify the limit of the exact monotone scheme. Stability makes the upper and lower half-relaxed limits finite. At a strict local maximum of the upper limit minus a smooth test function, select nearby maximizing grid points. Monotonicity compares the neighboring scheme values to that test function plus the vanishing contact offset. Local consistency then gives the viscosity subsolution inequality. The corresponding minima give the supersolution inequality for the lower limit. Consistent boundary data and the assumed barriers give the required boundary inequalities. Comparison orders the upper limit below the lower limit. The reverse order follows by definition, so the two limits coincide with the unique continuous viscosity solution. The vanishing distance between $v_h$ and $v_h^*$ transfers that convergence to the neural arrays. The discrete policy-value conclusion follows from the third bound. Interpreting a discrete policy as a continuous-time implementation additionally requires convergence of its controlled interpolation and payoff; the theorem does not infer this merely from convergence of its state arrays. This is the usual monotone-scheme selection argument, with the neural approximation error made explicit; see \citet{barles1991,crandall1992}.\qed

\subsubsection{Finite-horizon Bellman error accumulation}
Let $T_n$ and $T_n^a$ be the time-$n$ Bellman and fixed-policy maps, with Lipschitz constants $q_n$ (which may exceed one). For arrays $v_n$, define $\delta_n=\|v_n-T_n^av_{n+1}\|$ and $\epsilon_n=\|T_nv_{n+1}-T_n^av_{n+1}\|$. If the terminal error is $B_N$, backward induction gives
\[
 \|v_n-v_n^*\|\leq\left(\prod_{j=n}^{N-1}q_j\right)B_N
 +\sum_{k=n}^{N-1}\left(\prod_{j=n}^{k-1}q_j\right)(\delta_k+\epsilon_k).
\]
The policy-evaluation bound replaces $\delta_k+\epsilon_k$ by $\delta_k$. Their sum bounds regret. Empty products equal one. For the implicit Lipschitz recursive step, $q_n\leq(1-hL)^{-1}$; for additive exponential discount, $q_n=e^{-\rho h}$. Thus a per-step error $o(h)$ and consistent terminal data suffice on a fixed finite horizon, without asserting a stationary contraction in an inapplicable utility regime.

\subsubsection{Derivation of the temporal-self condition}
A self at $(t,x)$ replaces $a^*$ by a fixed feasible local action $a$ for duration $\epsilon$ and then returns to $a^*$. For each exponential component, the deviation value is
\[
 \E\left[\int_t^{t+\epsilon}e^{-\rho_j(s-t)}u(a)ds
     +e^{-\rho_j\epsilon}W_j(t+\epsilon,S_{t+\epsilon}^{a})\right].
\]
It\^o expansion gives its first-order increment relative to $W_j(t,x)$ as $\epsilon[u(a)+(W_j)_t+\LL^aW_j-\rho_jW_j]+o(\epsilon)$. Weighting the two components and subtracting their evaluation equations under $a^*$ leaves
\[
 \epsilon\{u(a)-u(a^*)+(\LL^a-\LL^{a^*})V\}+o(\epsilon).
\]
Nonpositive first-order gain for every feasible action is exactly \eqref{eq:spike}. This derivation assumes the smoothness and local integrability needed for the expansion; it is an equilibrium deviation condition, not an assertion of optimal commitment over a finite interval. The homothetic derivation follows by integrating $e^{-[\rho_j-b(r-m)]s}$ and applying the consumption first-order condition.\qed

% Reviewed source: revisions/2026-10-03-r9/manuscript/proofs.tex, line 1, commit f5021cefa71492babfcbfa580e0c984f59a9de26
\subsection{Proofs for the Guarded and Signed Map}\label{app:r9proofs}
\subsubsection{Proof of Proposition~\ref{prop:r9guard}}
At a replaced node the output equals the computed target. At any other node its distance from that target is at most $\tau_h$. The triangle inequality gives the one-step conclusion. For a frozen policy sequence, let $e_n$ be the sup-norm difference between the guarded and exact policy values at date $n$. Positivity of the transition weights and their mass bound imply $e_n\leq\tau_h+\eta_h+q_he_{n+1}$, with $e_N=0$ for exact terminal data. Backward induction gives the geometric sum. As $N=T/h$ and $q_h\leq1$, the bound is at most $T(0.1h+\eta_h/h)$ under the stated schedule. No property of the optimizer was used. The assertion concerns the guarded nodal values. An off-node neural interpolant, an improved policy, and a diffusion approximation each require their separate error account. In fixed-precision arithmetic, $\eta_h=o(h)$ is a maintained refinement condition, not a fact that can be inferred as $h$ tends to zero indefinitely.\hfill$\square$

\subsubsection{Proof of Theorem~\ref{thm:r9cover}}
On a regular box, the payoff restricted to a line segment is continuous and piecewise continuously differentiable, with finitely many interpolation and reflection knots, or is the limit of such segments. Integration of its almost-everywhere derivative yields
\[
 Q_w(a)-Q_w(c)=\int_0^1 \nabla Q_w(c+s(a-c))\cdot(a-c)\,ds.
\]
Every partial derivative is included in its interval hull. Maximizing the resulting interval products proves \eqref{eq:r9mean}. Taking the minimum with another valid upper enclosure preserves validity. If a partial derivative is strictly positive throughout the box, integration along that coordinate shows that replacing it by its upper endpoint cannot lower the payoff; the negative case is analogous. Repeating the argument proves face domination and the dimension statement.

A subdivision replaces a parent by children whose union contains it. A face reduction replaces a set by a subset with the same supremum. Removal occurs only after that supremum has a bound no larger than the incumbent plus tolerance. Feasible point evaluations can only increase the incumbent, so an earlier removed box stays bounded by the final incumbent plus tolerance. The remaining boxes keep their upper bounds without any success inference from an interrupted search. Induction over all operations proves \eqref{eq:r9coverbound}. At a stopping switch the line-segment continuity premise can fail; the implementation therefore uses the natural union enclosure and does not invoke the derivative argument there.

For the local quadratic observation, a $K$-Lipschitz gradient with a zero-gradient point in the box has norm at most $Ks$ throughout a box of diameter $s$. Integration over a segment of length at most $s$ bounds the variation by $Ks^2$. This conditional observation does not extend to an arbitrary switching box, nor does it assert that all nodes possess such a box.\hfill$\square$

\subsubsection{Derivative hulls for the preference economy}
Write $y=\log X$, and let $F$ be the continuation interpolant evaluated at a cubature transition. On an interior leg, its action derivatives are
\[
 \partial_c F=-hF_y,\qquad
 \partial_p F=F_y\{h(0.06-0.04p)+s_z\},\qquad
 \partial_\theta F=hF_u r_u,
\]
where $s_z$ is the stored portfolio shock coefficient and $r_u\in\{-1,1\}$ is the reflection slope; both slopes enter the hull at a fold. The flow derivatives are
\[
 h\exp\{(1-u)y-u\log c\},\qquad 0,\qquad -hk\theta.
\]
Outside a wealth stopping face, clipped wealth is constant and $F$ is the exact stopping payoff. Its $y$ derivative with respect to the transition is zero, while
\[
 F_u=\frac{\exp\{(1-u)y\}\{1-(1-u)y\}}{(1-u)^2}.
\]
Every crossed bilinear cell contributes its one-sided derivative range, restricted by the other coordinate's interval. Interval evaluation includes these ranges and the flow terms. A leg whose wealth interval meets both sides of a stopping switch disables the signed bound for that box. Point and interval evaluations share the stated nodal primitives but are computed independently of the actor's derivative graph.

\subsubsection{Proof of Proposition~\ref{prop:r9capital}}
Let $\overline m=d^{-1}\sum_i m_i$. Jensen's inequality gives $d^{-1}\sum_i\log m_i\leq\log\overline m$. Moreover, $d^{-1}\sum_i\tanh((BY)_i)\leq1$. The mean-state drift is consequently no larger than $b_+-\overline m$. The dispersion penalty at the terminal date is nonnegative and can be dropped for an upper bound. Fubini's theorem applied to the integrated expected mean state and terminal mean state yields the kernel $w(t)$. The Brownian terms have zero mean; bounded controls and bounded nonlinear drift ensure the required integrability. Thus every admissible policy has expected payoff at most
\[
 \int_0^T e^{-\rho t}\max_{m\in[\underline m,\overline m]}
 \{\log m-\zeta m^2/2+w(t)(b_+-m)\}\,dt.
\]
The integrand is strictly concave in $m$. Its derivative is $1/m-\zeta m-w(t)$, whose positive root gives $m_0(t)$ before clipping. This proves the upper bound.

For the feasible common schedule, the expected mean-state drift is at least $b_--m_0(t)$. Define the normalized dispersion norm $\|Z\|_{2,d}^2=\mathbb E[d^{-1}\sum_i Z_i^2]$. Centred common controls and common noise vanish from the dispersion equation. For numbers in $[-1,1]$, their cross-sectional variance is at most one. Hence Minkowski's inequality bounds the norm of the integrated centred nonlinear drift by $\kappa T$. The centred idiosyncratic Brownian part has norm $\sigma_I\sqrt{(1-1/d)T}$. A second application of Minkowski and squaring give the upper bound $D_d$ on expected terminal dispersion. Subtracting $\chi e^{-\rho T}D_d$ from the lower mean-payoff calculation proves $J(m_0)\geq L_d$. Feasibility gives $J(m_0)\leq V^*$. The derivation never restricts the dense matrix $B$ beyond the stated bounded-$\tanh$ form.\hfill$\square$

\subsubsection{Compensation and arithmetic scope}
For each realized state and positive consumption, CRRA felicity is increasing in consumption, including when $u>1$ and its level is negative. The financed supplement leaves the state law and every other payoff term unchanged. Its expected payoff is therefore nondecreasing in $e$. Positivity of the nodal transition operator propagates a lower interval payoff at each date. If its lower endpoint exceeds the optimal upper envelope at every active starting node and date, the corresponding financed supplement compensates the loss throughout that nodal economy. Bisection searches for a sufficient supplement; only the final verified inequality provides the stated upper bound.

All implemented interval conclusions use the inherited binary64 execution contract, with outward elementary operations and separately bounded transcendental evaluations. They concern exact real interpretations of stored primitive coefficients and grid nodes. They are not a formal verification of the interpreter, hardware, or all library internals. In the capital bound, square-root seeds are accepted only after interval squaring verifies both endpoint inequalities, and the time integral is bounded by interval rectangles rather than by an adaptive quadrature tolerance. Random stress tests are regression evidence, not substitutes for these inclusion arguments.
% Reviewed source: revisions/2026-10-04-r14/manuscript/proofs_r12.tex, line 1, commit f5021cefa71492babfcbfa580e0c984f59a9de26
\subsection{Additional Proofs and Numerical Contracts}\label{app:r12proofs}
\subsubsection{Costate error and action quality}
Let $\widehat m$ maximize $H_{\widehat q}$ on the stated convex feasible set, and let $e=q-\widehat q$. Strong concavity and the variational first-order condition at $\widehat m$ imply
\[
 H_{\widehat q}(\widehat m)-H_{\widehat q}(m)
       \geq \frac\mu2\|m-\widehat m\|_d^2
\]
for every feasible $m$, including boundary maximizers. Therefore
$\|\widetilde m-\widehat m\|_d\leq\sqrt{2\delta/\mu}$. Decomposing the true Hamiltonian difference yields
\begin{align*}
 H_q(m_q)-H_q(\widetilde m)
 &\leq \delta-\frac\mu2\|m_q-\widehat m\|_d^2
           +\langle e,\widetilde m-m_q\rangle_d\\
 &\leq \delta+\|e\|_d\sqrt{2\delta/\mu}+\frac{\|e\|_d^2}{2\mu}.
\end{align*}
The last step completes the square in $\|m_q-\widehat m\|_d$. This is \eqref{eq:r12actionloss}. The Hessian of $H_q$ is the negative diagonal matrix with entries $1/(dm_i^2)$ minus the positive semidefinite mean-penalty matrix. It is bounded above by $-\mu I/d$, proving the strong-concavity premise.

For \eqref{eq:r12projection}, write $Z-g(X)=[Z-\E(Z\mid X)]+[\E(Z\mid X)-g(X)]$. The cross term has zero expectation by conditioning on $X$. Integrability justifies this orthogonal decomposition. If the conditional Euler costate differs from the continuous policy costate by $b(X)$, that term remains an additional bias; it is not removed by the projection identity. Neither the sampled costate loss nor a successful optimizer return is a uniform upper bound on this bias.

\subsubsection{Cancellation, clipping and finite-population inference}
For each candidate, the exact payoff identity used in Proposition~\ref{prop:r11transfer} consists of a policy production term, its consumption term, its terminal dispersion, and the corresponding reference terms. When two candidates use the identical reference path, subtraction cancels every exact reference term. Comparing the two policy terms with their numerical approximations therefore requires the two policy transfer errors and the statistic-roundoff accounts, but no reference-state transfer error. The implementation retains conservative original statistic budgets even where their reference quadrature components also cancel.

For any real $x,y$ and positive thresholds $b_a,b_b$,
\[
 (|x-y|-b_a-b_b)_+
       \leq (|x|-b_a)_++(|y|-b_b)_+.
\]
Taking expectations proves the summed clipping-tail allowance without requiring independence between policies. Applying the empirical Bernstein inequality to the clipped difference and then to its negative gives the two one-sided bounds. A union bound covers all declared tests. Our allocation counts both directions and every method, profile design and auxiliary evaluation; it does not spend new confidence mass for deterministic transformations such as \eqref{eq:r12fee}.

For the finite population, draw $I$ with the fixed probabilities $p_j$, independently of each future innovation path. The pair $(I,W)$ is then independent across paths. Conditional transfer gives
\[
 \left|\E[X_{a,h}]-\sum_jp_j\{J_{y_j}(a)-J_{y_j}(m_0)\}\right|
       \leq\sum_jp_j E_{a,h}(y_j).
\]
The clipping threshold is the maximum of the profile-specific thresholds; its expectation-tail allowance is bounded by the maximum of their allowances. The prescribed implementation uses equal probabilities over nine fixed profiles. Both averaging and the final subtraction are interval-enclosed. Random empirical profile frequencies are not substituted for the population weights when calculating deterministic error budgets. They are part of the sampling variation of the final statistic.

For the paired method difference, the same profile is assigned to both policies on each path. Common random numbers across different methods or training seeds do not invalidate a finite-family union bound. They also do not create independent replications of the optimizer. The ten fixed training streams are reported individually and descriptively; no confidence assertion about an unspecified distribution of future hyperparameter choices is inferred.

\subsubsection{Observed-history implementation}
Under the observation assumptions, $M$ is a continuous martingale with bracket $\langle M\rangle_t=Qt$. Let $A=\Sigma^\top Q^{-1}$. The matrix $A\Sigma$ is the orthogonal projector onto the row space of $\Sigma$, and $nn^\top$ is the projector onto its orthogonal complement. Independence of the private Brownian motion gives
\[
 \langle AM+nB^0\rangle_t
       =(AQA^\top+nn^\top)t=I_{d+1}t.
\]
The multivariate L\'evy characterization proves the Brownian assertion. Also $\Sigma A=I_d$ and $\Sigma n=0$, so $\Sigma\widetilde W=M$. On each decision cell, the action is determined at its left endpoint by the stored internal state and past recovered increments. With that held action, the physical state equation has a unique strong solution because the production drift is globally Lipschitz and bounded. Induction over the cells establishes uniqueness in law for the joint physical and internal process. Driven by the Brownian motion just constructed, it therefore has the same law as the controller originally defined using directly supplied innovations. Applying the specified componentwise clipping to the recovered vector preserves this distributional equivalence.

The proposition uses only private randomization and observed state history; it does not reveal the original null-space economic driver. Randomized controls are admissible in the enlarged filtration. Since the original deterministic performance envelope holds against all such adapted controls and independent randomization cannot improve an unconstrained expected-payoff supremum beyond the best deterministic randomization realization, the comparison with the original optimum is preserved. The result is not an implementability theorem for discrete noisy state measurements. The exact drift integral is stated explicitly, and errors in approximating it must be treated separately.

\subsubsection{Independent scalar HJB approximation}
For $d=1$, the terminal dispersion penalty vanishes but production remains nonlinear in log capital. The reference solves the same primitive action-box problem, not a Riccati or homothetic reduction. Put $f(y)=\kappa\tanh(By)$ and $\nu=(\sigma_I^2+\sigma_C^2)/2$. On a uniform grid of width $\Delta y$, define the positive rates
\[
 \ell_i(m)=\frac\nu{\Delta y^2}+\frac{(m-c-f(y_i))_+}{\Delta y},
 \qquad r_i(m)=\frac\nu{\Delta y^2}+\frac{(c+f(y_i)-m)_+}{\Delta y}.
\]
The implicit Bellman step solves
\[
 \frac{v_i-v_i^+}{h}+\rho v_i
 =\max_m\{\log m+y_i-\zeta m^2/2
       +\ell_i(m)(v_{i-1}-v_i)+r_i(m)(v_{i+1}-v_i)\}.
\]
For a fixed action vector the coefficient matrix is strictly diagonally dominant with positive diagonal and nonpositive off-diagonal entries. Its inverse is nonnegative. Howard updates maximize over each drift-sign interval; on either interval the positive root of $1/m-\zeta m=p$ is clipped to its feasible endpoints, and the two candidates are compared. This is an independent monotone matrix computation, not a trained critic used as its own reference.

At the two outer boundaries, the prescribed affine asymptotes replace production by its limiting saturation value. If $w(t)$ is the coefficient on initial log capital and $m_0(t)$ the common schedule, their intercept is the numerical integral of the corresponding deterministic discounted flow. The reference reports the grid, time step, boundary distance, linear-system residual, policy-iteration convergence and domain sensitivity. Fitted actors are evaluated as Markov feedbacks on this finite grid. These comparisons do not conflate that controller with the innovation-driven implementation. Grid differences and asymptotic boundary values are not asserted to be rigorous continuous-domain error enclosures.

% Reviewed source: revisions/2026-10-04-r14/manuscript/mechanism_welfare.tex, line 1, commit f5021cefa71492babfcbfa580e0c984f59a9de26
\subsubsection{From evaluation error to economic improvement}
\label{sec:r14costate}
Proposition~\ref{prop:r12costate} controls a decision at a given state.
The distribution of states at which that decision is used matters for
the economic conclusion. We make this link explicit. In this subsection,
$\pi$ is a fixed Markov policy with a sufficiently regular policy value
$V^\pi$ to justify It\^o's formula and the expectations below. Specifically, $V^\pi\in C^{1,2}([0,T)\times\mathbb R^d)$ is continuous at the terminal date with $V^\pi(T,\cdot)=g$. Along each candidate considered, the drift and terminal terms in the discounted It\^o identity are integrable and its stochastic integral is a zero-mean martingale. The condition
\[
 \E\int_0^T e^{-2\rho t}
 \|\Sigma^\top D_yV^\pi(t,Y_t^a)\|_2^2\,dt<\infty
\]
is sufficient for the latter. The parameters of $\pi$ are frozen. Let $a$ be any admissible candidate, including
a history-dependent candidate, and let $Y^a$ be its physical state.
Set $q^\pi=dD_yV^\pi$ and define
\[
 g^\pi(t,y)=\max_{m\in\mathcal A}H_{q^\pi(t,y)}(m)
                 -H_{q^\pi(t,y)}(\pi(t,y))\geq0,
\]
where $\mathcal A$ is the same convex action set for both maximizations
and both policies. One may always take the primitive action box.
For a costate predictor $\widehat q$, let $\delta_t$ bound the
$H_{\widehat q(t,Y_t^a)}$ maximization gap of $a_t$ over that set.
Write $\overline q_h(t,y)$ for the conditional mean of the frozen-policy
rollout costate at mesh $h$. Define, under the candidate's state law,
\begin{align*}
 \mathcal G_a&=\E\int_0^T e^{-\rho t}g^\pi(t,Y_t^a)\,dt,\\
 \mathcal D_a&=\E\int_0^T e^{-\rho t}\delta_t\,dt,\\
 \mathcal E_{a,h}^2&=\E\int_0^T e^{-\rho t}
      \|\widehat q(t,Y_t^a)-\overline q_h(t,Y_t^a)\|_d^2\,dt,\\
 \mathcal B_{a,h}^2&=\E\int_0^T e^{-\rho t}
      \|\overline q_h(t,Y_t^a)-q^\pi(t,Y_t^a)\|_d^2\,dt.
\end{align*}
The quantity $\mathcal E_{a,h}$ concerns regression accuracy;
$\mathcal B_{a,h}$ retains the rollout discretization bias.

\begin{proposition}[Evaluation, action quality, and payoff]
\label{prop:r14economiccostate}
Suppose the displayed quantities are finite and the consumption
Hamiltonian is $\mu$-strongly concave on $\mathcal A$. Then
\begin{equation}\label{eq:r14economiccostate}
 J_y(a)-J_y(\pi)\geq
 \mathcal G_a-
 \left\{\sqrt{\mathcal D_a}
       +\frac{\mathcal E_{a,h}+\mathcal B_{a,h}}{\sqrt{2\mu}}
 \right\}^{\!2}.
\end{equation}
In particular, an upper bound on the action-loss term and a lower bound
on $\mathcal G_a$ establish a sufficient condition for policy
improvement. The proposition applies to any costate estimator; a
method comparison still requires the independently evaluated payoffs
and their associated work.
\end{proposition}

\begin{proof}
The fixed-policy equation for $V^\pi$, evaluated along $Y^a$, gives
the performance-difference identity
\begin{equation}\label{eq:r14performanceidentity}
 J_y(a)-J_y(\pi)=\E\int_0^T e^{-\rho t}
  \{H_{q^\pi(t,Y_t^a)}(a_t)
       -H_{q^\pi(t,Y_t^a)}(\pi(t,Y_t^a))\}\,dt.
\end{equation}
The diffusion and the state-dependent part of felicity are common to
the two actions, so their differences cancel. The terminal terms agree.
Localization followed by the assumed integrability justifies the
martingale expectation. Subtract the true greedy Hamiltonian inside
the integrand and use Proposition~\ref{prop:r12costate}. Minkowski's
inequality for the measure $e^{-\rho t}dt\,d\Pr$, first for
$\sqrt{\delta_t}$ and the costate norm, and then for the regression
and bias components of the costate, gives
\[
 \E\int_0^T e^{-\rho t}
 \left(\sqrt{\delta_t}+
     \frac{\|q^\pi-\widehat q\|_d}{\sqrt{2\mu}}\right)^2dt
 \leq\left\{\sqrt{\mathcal D_a}+
       \frac{\mathcal E_{a,h}+\mathcal B_{a,h}}{\sqrt{2\mu}}\right\}^2.
\]
Combining these statements proves the result.
\end{proof}

\paragraph*{Independent costate diagnostics.}
Condition on a fitted policy, its critic, and a held-out state $X$.
Let $\overline Z_1,\overline Z_2$ be means from independent rollout
banks with the same conditional mean $\overline q_h(X)$. Conditional
independence and zero conditional mean of the two errors imply
\[
 \E[\langle\widehat q(X)-\overline Z_1,
                  \widehat q(X)-\overline Z_2\rangle_d\mid X]
       =\|\widehat q(X)-\overline q_h(X)\|_d^2.
\]
With $r$ independent replicates per bank and conditional single-rollout
variance $v_h(X)=\E[\|Z-\overline q_h(X)\|_d^2\mid X]$,
\[
 \frac12\E[\|\overline Z_1-\overline Z_2\|_d^2\mid X]
       =\frac{v_h(X)}r.
\]
The first sample average is not constrained to be positive. Both
identities refer to the actual held-out state distribution. They
estimate the quantities in~\eqref{eq:r14economiccostate} only when
that distribution is the candidate's discounted state occupation law,
or when a justified change-of-measure account is supplied. Nested-grid
differences diagnose time-step sensitivity; without a remainder bound
they do not bound $\mathcal B_{a,h}$.

\paragraph*{A continuous-time costate error account.}
There is also a direct sufficient condition for controlling the true
costate, independent of a rollout bias estimate. For a smooth candidate
value $v$, define its fixed-policy residual and terminal error by
\[
 r^\pi[v]=-v_t-\mathcal L^\pi v-u(\cdot,\pi)+\rho v,
 \qquad b_v=v(T,\cdot)-g.
\]
Suppose $b^\pi=c\bm1+f-\pi$ is continuously differentiable with
$\|D_yb^\pi\|_2\leq L_b$, and suppose differentiating the following
Feynman--Kac expectations is justified. If
$\|dD_yb_v\|_d\leq C_T$ and
$\|dD_yr^\pi[v](s,\cdot)\|_d\leq C_r(s)$ everywhere, then
\begin{equation}\label{eq:r14costateresidual}
 \|dD_yv(t,y)-q^\pi(t,y)\|_d
 \leq e^{(L_b-\rho)(T-t)}C_T
       +\int_t^T e^{(L_b-\rho)(s-t)}C_r(s)\,ds.
\end{equation}
To see this, write $v-V^\pi$ as the expected discounted terminal error
plus the expected discounted residual integral. The state Jacobian
$J_{t,s}$ satisfies
$\dot J_{t,s}=D_yb^\pi(s,Y_s)J_{t,s}$, $J_{t,t}=I$, and hence
$\|J_{t,s}\|_2\leq e^{L_b(s-t)}$. Differentiate the expectation and
apply that bound. A hard terminal condition gives $C_T=0$.
Equation~\eqref{eq:r14costateresidual} is a sufficient mathematical
account for a smooth continuous policy. A sampled residual, a rollout
grid difference, or the derivatives of a different deployed
implementation do not supply its global hypotheses.
